\movedsub{The same probe with entities instead of numbers}
\label{sec:entity}

A numeric substitution is one error type, and the pattern above could be specific to it. The same construction runs on named entities: take a grounded entity in the answer, replace it with a donor entity that the context does not contain, hold everything else fixed. Cases are built once and every evaluator sees the same ones. Table~\ref{tab:entity} reports them.

The verbosity effect survives the change of error type. Under the entailment aggregate on Natural Questions a falsified entity moves the score by 0.082 for the most verbose generator and 0.403 for the tersest, against 0.054 and 0.451 for numbers. Detection at the fixed gate is 9.1\% against 48.8\%, against 5.1\% and 56.6\%. The gap narrows slightly and does not close.

Ordering by generator holds in five of the six dataset-evaluator combinations. The exception is claim-min on Natural Questions, where the verbose generator is scored \emph{more} sensitively than the mid-length one (0.163 against 0.112, detection 25.9\% against 17.9\%). Assertion count predicts the ordering well but it does not fix it.

Entity and numeric substitution are not interchangeable, and neither is uniformly the harder case. Both agree closely for the tersest generator, within 0.05 in every cell. They diverge for the other two, and in opposite directions: on HotpotQA the entailment aggregate responds roughly twice as strongly to a falsified entity as to a falsified number for the verbose generator (0.178 against 0.085), while on Natural Questions the mid-length generator shows the reverse and by a wider factor (0.088 against 0.261). A claim that an evaluator does or does not catch fabrications is therefore a claim about a particular kind of fabrication, and the two error types we tested do not license extrapolation to each other, let alone to unsupported causal claims or contradicted negations, which we did not test.

\begin{table*}[t]
\centering
\small
\begin{tabular}{llrrrrrr}
\toprule
 & & \multicolumn{2}{c}{\textbf{Claude Haiku 4.5}} & \multicolumn{2}{c}{\textbf{Qwen2.5-7B}} & \multicolumn{2}{c}{\textbf{GPT-4o-mini}} \\
\cmidrule(lr){3-4}\cmidrule(lr){5-6}\cmidrule(lr){7-8}
\textbf{Dataset} & \textbf{Evaluator} & $\Delta$ & Det. & $\Delta$ & Det. & $\Delta$ & Det. \\
\midrule
NQ & NLI-max & 0.082 & 9.1\% & 0.088 & 10.8\% & 0.403 & 48.8\% \\
 & Claim-min & 0.163 & 25.9\% & 0.112 & 17.9\% & 0.431 & 53.4\% \\
 & AlignScore & 0.232 & 21.9\% & 0.290 & 33.1\% & 0.673 & 77.5\% \\
\midrule
HotpotQA & NLI-max & 0.178 & 28.6\% & 0.267 & 54.1\% & 0.412 & 75.1\% \\
 & Claim-min & 0.132 & 38.2\% & 0.262 & 67.3\% & 0.412 & 75.9\% \\
 & AlignScore & 0.321 & 47.6\% & 0.463 & 69.4\% & 0.668 & 90.6\% \\
\bottomrule
\end{tabular}
\caption{Response to a falsified \emph{entity}, on the same answers and the same fixed 0.5 gate as Table~\ref{tab:2}. $\Delta$ is the mean drop in score. n = 783, 740 and 725 on NQ and 800, 771 and 762 on HotpotQA; detection is computed over the cases scoring above the gate before falsification. The generator ordering of Table~\ref{tab:2} is reproduced in five of six rows, the exception being claim-min on NQ.}
\label{tab:entity}
\end{table*}
